% Part H supplementary note.  This is a TeX fragment; the main document may
% \input it at the desired position outside, or in a suitable column of, paracol.

\section*{スカートについて}

\begin{japanese}
スカート（\textgerman{Skat}）は、三人で行うドイツの伝統的な札遊びである。
三十二枚の札を用い、一人に十枚ずつ配り、残る二枚を伏せておく。
この二枚もまたスカートと呼ばれる。一局ごとに一人が\emph{独り方}となり、
ほかの二人を相手に札を取って、その札点を競う。

本書の会話が前提とするのは、今日一般的な数字による競りよりも古い
「組による競り」である。競る者は、グリューン・ソロ、ヌル、グランなど、
しようとする勝負の種類を順に申し出る。ほかの者がその競りを受けずに
パスすれば、最後まで残った者が独り方となる。グランでは四枚の
ヴェンツェルだけが切り札となり、そのほかの札はそれぞれの組に従う。

ここで用いられる札は、アイヒェル（団栗）、グリューン（葉）、
ロート（赤）、シェレン（鈴）の四組からなるドイツ式の札である。
ヴェンツェルは各組のウンターで、通常の勝負では最も強い四枚の切り札となる。
ダウスは各組で十一点を持つ高札である。最初に出された組の札を持つ者は、
同じ組の札を出さねばならない。三人の札のうち最も強いものを出した者が
その一巡の札を取り、次の札を最初に出す。

独り方は、取った札とスカートを合わせて六十一点以上を得れば勝つ。
一方の側が三十点以下にとどまることをシュナイダーという。
したがって「シュナイダーを免れた」とは、少なくとも三十一点を得て、
この重い負け方だけは避けたという意味である。

第二の会話の「四枚持ちのグラン、六十点」は、最上位から続く四枚の
ヴェンツェルを独り方が持つものとして数えた言い方である。
グランの基礎値を十二とする当時の地方的な勘定法では、
四枚持ちに勝負そのものの一を加え、十二の五倍で六十点となる。
スカートの規則と勘定法は時代や土地によって異なったため、
ここでは現行規則に改めず、会話に示された古い遊び方をそのまま読むのがよい。
\end{japanese}
